\section{RQ5: Declared versus Deployed on Federal Systems}
\label{sec:fed}

RQ2 measured what federal endpoints negotiate and RQ3 what open-source code
contains. Neither shows both sides of one system. Federal agencies publish much
of their code, which lets us put the two side by side for the same service.

\subsection{Method}

\paragraph{Linking code to endpoints.} GitHub maintains a public registry of
government accounts; we took the \data{fed-orgs} organizations listed under
U.S.\ Federal (registry commit \texttt{\data{fed-registry-commit}}) and listed
their public repositories on \data{fed-date}: \data{fed-repos} repositories
(\data{fed-orgs-unlistable} organizations could not be listed). A repository
is linked to an endpoint when its homepage field names exactly one of the hosts
probed in RQ2. That gives \data{fed-links} repositories for
\data{fed-link-endpoints} endpoints, of which \data{fed-current} for
\data{fed-current-endpoints} endpoints are current: not archived, not a fork,
and linked to an endpoint that speaks TLS.

A homepage field often names the agency's site for code that has nothing to do
with running it, so we separate the links. A link is \emph{strong} when the
repository is the site's own source, shown by a GitHub Pages \texttt{CNAME}
file naming the host or by a repository name that contains the whole domain;
\data{fed-strong} links are strong and \data{fed-weak} weak. The rule is
conservative: site repositories named differently count as weak.

\paragraph{Reconciling.} Each repository was cloned and scanned as in RQ3, with
the endpoint's RQ2 records merged into the same component (\S\ref{sec:design}),
so that reconciliation can compare what the code declares with what the
endpoint negotiates. No endpoint was probed again.
\data{fed-scanned} repositories were scanned; \data{fed-clonefail} failed to
clone.

\subsection{Results}

\begin{table*}[t]
\centering\small
% Generated by paper/gen.py from paper/data. Do not edit.
\begin{tabular}{llllrr}
\toprule
Repository & Site & Deployed & Hosting AS & PQ & Vuln. \\
\midrule
\texttt{GSA/councils.gov} & councils.gov & classical & MetTel & 0 & 0 \\
\texttt{GSA/data.gov} & data.gov & hybrid & Amazon (AMAZON-02) & 0 & 0 \\
\texttt{cisagov/get.gov} & get.gov & hybrid & Cloudflare & 0 & 0 \\
\texttt{cisagov/manage.get.gov} & get.gov & hybrid & Cloudflare & 2 & 14 \\
\texttt{cfpb/consumerfinance.gov} & www.consumerfinance.gov & hybrid & Akamai & 0 & 0 \\
\texttt{usdoj/foia.gov} & www.foia.gov & hybrid & Akamai & 0 & 0 \\
\texttt{GSA/Section508.gov} & www.section508.gov & hybrid & Amazon (AMAZON-02) & 0 & 0 \\
\texttt{usds/website} & www.usds.gov & hybrid & Fastly & 0 & 0 \\
\bottomrule
\end{tabular}

\caption{The \data{fed-strong-scanned} strong links: repositories that are a
federal site's source, the group the site negotiated in RQ2, its hosting AS,
and the post-quantum (PQ) and quantum-vulnerable (Vuln.) findings in the
repository.}
\label{tab:fed-strong}
\end{table*}

\paragraph{The gap runs the other way.} \tool{} was built to catch code that
declares post-quantum support in front of an endpoint that deploys classical
key exchange. On federal systems that never happens: no repository produced a
declared-PQC-deployed-classical finding, and none produced a
configuration-overrides-library finding. Only \data{fed-declares} of the
\data{fed-scanned} repositories declares post-quantum capability at all,
\data{fed-declares-list}, through its cryptography library. Yet
\data{fed-hybrid-endpoints} of the \data{fed-endpoints} linked endpoints
(\data{fed-hybrid-endpoints-share}) negotiate a hybrid group, behind
\data{fed-hybrid-repos} of the repositories. Deployment is ahead of
declaration.

\paragraph{Protection comes from the platform.} \data{fed-nocrypto} of the
\data{fed-scanned} repositories (\data{fed-nocrypto-share}) contain no
public-key cryptography that \tool{} detects, post-quantum or vulnerable; only
\data{fed-vulnerable-repos} contain any quantum-vulnerable findings. Every
hybrid endpoint is served from one of \data{fed-hybrid-providers} providers:
\data{fed-hybrid-provider-list}. Table~\ref{tab:fed-strong} shows the pattern
for the strong links. The Consumer Financial Protection Bureau's site, the
Justice Department's FOIA portal and the U.S.\ Digital Service's site all
negotiate X25519MLKEM768, and their source contains no public-key cryptography
that \tool{} detects. Their addresses belong to content-delivery networks,
which terminate TLS in front of the application. The one strong
link that negotiates a classical group, \texttt{councils.gov}, is hosted by a
telecommunications carrier, and its source is equally silent.

\paragraph{What follows for inventories.} M-26-15 asks agencies to inventory
their cryptography, and a natural first step is the code they maintain. For
public web services that step finds nothing to migrate and cannot tell
whether the service is quantum-safe: in our sample the answer depends on a
platform the code never mentions. Measuring the deployed endpoint, as the
network layer does, is the only one of \tool{}'s layers that answers it.

\paragraph{Limits.} The strong sample is small, \data{fed-strong} pairs, and
the weak links include repositories that only mention the agency's site. The
homepage field is self-reported; repositories that build a site without naming
it are missed. Applications behind these sites may use cryptography in services
whose code is not public. The endpoint data are RQ2's, from one vantage point
on \data{rqtwo-date}.
